% Appendix 7: Linear Dependence and Independence
% VIC Specialist Mathematics — Units 3–4
\begingroup
\small

\noindent\fbox{%
\begin{minipage}{0.97\textwidth}
\textit{\textbf{Enrichment only.}
Linear dependence and independence are part of
\textbf{VIC Specialist Mathematics (Units~3--4)}.
The NESA syllabus uses these concepts informally in geometric proofs
(e.g.\ showing coplanarity via linear combinations);
this appendix provides the formal definitions.}
\end{minipage}}

\vspace{0.8em}

% ---- Key concepts ----
\noindent\textbf{Key concepts.}

A \textbf{linear combination} of vectors is an expression formed by
multiplying vectors by scalars and adding them.
For example,
\[
  \mathbf{w} = 2\mathbf{a} - 3\mathbf{b}
\]
is a linear combination of $\mathbf{a}$ and $\mathbf{b}$.

\vspace{0.4em}
Vectors $\mathbf{v}_1, \ldots, \mathbf{v}_n$ are
\textbf{linearly independent} if
\[
  \boxed{
  c_1 \mathbf{v}_1 + \cdots + c_n \mathbf{v}_n = \mathbf{0}
  }
\]
is possible \textbf{only} when
\[
  c_1 = c_2 = \cdots = c_n = 0.
\]

Otherwise, they are \textbf{linearly dependent}.

\vspace{0.4em}
Informally, dependence means at least one vector can be constructed
from the others: it does not provide a genuinely new direction.

% ---- Worked example ----
\vspace{0.8em}
\noindent\textbf{Worked example.}

\noindent Consider the three vectors
\[
  \mathbf{a} = (1, 0, 1), \qquad
  \mathbf{b} = (0, 2, 0), \qquad
  \mathbf{c} = (1, 2, 1).
\]
Determine whether they are linearly independent.

\vspace{0.3em}
\noindent\textit{Solution.}
Observe that
\[
  \mathbf{a} + \mathbf{b} = (1, 2, 1) = \mathbf{c}.
\]
Therefore,
\[
  \boxed{\mathbf{a} + \mathbf{b} - \mathbf{c} = \mathbf{0}}.
\]
The coefficients $1, 1, -1$ are not all zero.
Hence, the three vectors are \textbf{linearly dependent}.

\vspace{0.4em}
\noindent\textit{Geometric interpretation.}
All three vectors lie in the plane $x - z = 0$ through the origin.

However, $\mathbf{a}$ and $\mathbf{b}$ themselves \emph{are}
independent. Indeed,
\[
  s\,\mathbf{a} + t\,\mathbf{b}
  = (s,\, 2t,\, s) = \mathbf{0}
\]
implies $s = t = 0$.

Thus, two independent vectors can describe the whole plane, while
the third adds no new direction.

% ---- Key takeaways ----
\vspace{0.8em}
\noindent\textbf{Key takeaways.}
\begin{itemize}[leftmargin=*,itemsep=2pt]
  \item Linear dependence means some vector information is redundant.
  \item Two non-zero vectors are independent precisely when they are
        \textbf{not parallel}.
  \item In three-dimensional space, three independent vectors can
        describe every vector through linear combinations.
        Any set of four or more 3D vectors is necessarily dependent.
\end{itemize}

\vspace{0.4em}
\noindent\textit{Beyond school --- University connection.}
Linear independence leads to the concept of a \emph{basis} in linear
algebra. It is fundamental to matrix rank, solving simultaneous
equations, computer graphics and data science.

\endgroup
